% =====================================================================
% Geographic encoding in premodern enumerations:
% a five-corpus quantitative test (RV 10.75.5, Mahajanapadas,
% Avesta Fargard 1, Bisutun dahyava, Genesis 10).
%
% Compile with lualatex via the centralized harness:
%   make Nadīstuti-sūkta
% (see /Users/akulkarn/Desktop/Code/Ajinkya/Papers/Makefile)
% =====================================================================

\documentclass[11pt]{article}

\usepackage[a4paper, margin=1in]{geometry}
\usepackage{fontspec}
\setmainfont{Latin Modern Roman}
\usepackage{microtype}

\usepackage{amsmath, amssymb, amsfonts}
\usepackage{booktabs}
\usepackage{tabularx}
\usepackage{array}
\usepackage{graphicx}
\graphicspath{{figures/}}

\usepackage[round, authoryear]{natbib}

\usepackage{xcolor}
\definecolor{linkcolor}{rgb}{0.55,0.14,0.20}
\usepackage[colorlinks=true,
            linkcolor=linkcolor,
            citecolor=linkcolor,
            urlcolor=linkcolor]{hyperref}
\usepackage{url}

\usepackage{fancyhdr}
\pagestyle{fancy}
\fancyhf{}
\fancyhead[L]{\small Geographic encoding in premodern enumerations}
\fancyhead[R]{\small Kulkarni 2026}
\fancyfoot[C]{\thepage}
\renewcommand{\headrulewidth}{0.4pt}

\usepackage{titlesec}
\titleformat{\section}{\large\bfseries\sffamily}{\thesection}{0.5em}{}
\titleformat{\subsection}{\normalsize\bfseries\sffamily}{\thesubsection}{0.5em}{}
\titlespacing*{\section}{0pt}{8pt}{4pt}
\titlespacing*{\subsection}{0pt}{6pt}{2pt}

\usepackage{caption}
\captionsetup{font=small, labelfont=bf, format=plain, skip=4pt}

\title{\bfseries Geographic encoding in premodern enumerations:\\
       a five-corpus quantitative test\\[2pt]
       \large\mdseries Ṛgveda 10.75.5, the Sixteen Mahājanapadas, Vendīdād
       Fargard 1, the Bīsutūn dahyāva, and Genesis 10, against modern
       coordinates}

\author{
  Ajinkya Kulkarni\thanks{Correspondence:
    \href{mailto:ajinkya.kulkarni@uah.edu}{ajinkya.kulkarni@uah.edu}.
    Independent scholarly work conducted outside the author's employment;
    views expressed are the author's alone.}\\
  University of Alabama in Huntsville
}

\date{May 2026}

\begin{document}
\maketitle
\thispagestyle{fancy}

\begin{abstract}
\noindent
Premodern textual traditions across Eurasia preserve enumerative lists
of rivers, lands, and peoples. The qualitative observation that such
lists sometimes ``trace a geographic sweep'' has been made informally
since the nineteenth century, but has rarely been tested with
quantitative method. We test the hypothesis that five canonical
enumerations encode monotonic geographic order in the sequence of
items they name: the ten-river enumeration of Ṛgveda 10.75.5
(Nadīstuti-sūkta), the sixteen Mahājanapadas of the Anguttara
Nikāya, the sixteen lands of Vendīdād Fargard 1, the twenty-three
subject lands (dahyāva) of Darius I's Bīsutūn inscription, and
the Table of Nations in Genesis 10. For each corpus we compute Kendall's $\tau$
between sequence position and modern longitude (and latitude) of the
identified items, using a permutation null distribution, a Procrustes
superimposition test, and a Bayes factor comparison between
geographic and uniform-random generative models. The two Indic
corpora (RV 10.75.5 and the Mahājanapadas) show strong monotonic
east-to-west ordering ($\tau = -0.911$ and $-0.783$ respectively,
both $p<10^{-4}$). Bīsutūn dahyāva shows moderate west-to-east
ordering with strong pairwise spatial clustering ($\tau = +0.47$,
$p = 0.001$; Mantel $r = +0.50$, $p<10^{-4}$). Vendīdād Fargard 1
shows no detectable geographic structure on any test. Genesis 10
shows only a weak monotonic signal, but significant pairwise spatial
coherence (Mantel $r=0.17$, $p=0.0004$): each of its three Noachic
sub-sections is latitudinally ordered, though the sub-sections run
in opposite directions and cancel at the surface. The sharpest
contrast is Bīsutūn against Vendīdād---both Iranian, both of
roughly the same date, naming much the same regions, yet the
administrative inscription is geographically ordered and the
cosmological text is not. Across the five, genre predicts geographic
ordering better than language family does. We make no claims about
the date, composition, or cultural identity of any of the five
traditions; the tests bear only on the texts as transmitted.
\medskip

\noindent\textbf{Keywords:} quantitative philology; rank correlation;
permutation test; Procrustes analysis; Bayes factor; comparative
textual analysis; historical geography.
\end{abstract}

% =====================================================================
\section{Introduction}\label{sec:intro}

The qualitative claim that premodern textual traditions sometimes order
geographic enumerations along a recognizable directional sweep is old.
\citet{macdonell-keith1912} note in passing that the river list of
\d{R}gveda~10.75.5 ``proceeds broadly from east to west.''
\citet{rhysdavids1903} observes a similar east-to-west tendency in
the Anguttara Nikāya's canonical list of sixteen mahājanapadas.
\citet{gnoli1980} and \citet{witzel2000homeland} discuss the Avesta's
Vendīdād Fargard 1 as a possible geographic itinerary across the
Iranian and Indo-Iranian world. \citet{westermann1984} and
\citet{wenham1987} treat the Hebrew Bible's Table of Nations
(Genesis~10) as encoding the post-deluge geographic distribution of
peoples ``known to the Israelites,'' implicitly with some geographic
structure.

All of these claims are informal. Each carries the same testable
content: that the order in which named entities appear is not
arbitrary with respect to where those entities sit on a map. We are
not aware of any study that has checked that content quantitatively,
let alone across several traditions at once.

We test five lists here, each well studied in its own field:
\d{R}gveda~10.75.5 (ten rivers; Vedic Sanskrit, late 2nd millennium
BCE); the sixteen Mahājanapadas of Anguttara Nikāya I.213 (Pali,
mid-1st millennium BCE); the sixteen lands of Vendīdād Fargard 1
(Avestan, mid-1st millennium BCE); the twenty-three subject lands of
Darius~I's Bīsutūn inscription (Old Persian, \textit{c.}\ 520 BCE);
and Genesis~10's Table of Nations (Biblical Hebrew, compiled in the
6th century BCE from older material). Two are Indic, two Iranian,
one Semitic; together they span close to a millennium. Each is
analysed with the same non-parametric methods.

\subsection{Scope and limits of the claim}\label{sec:scope}

The claim tested here is narrow. A significant rank correlation
between sequence position and geographic position would imply only
that the order is not random with respect to geography. It would not
imply any specific dating, any specific direction of cultural
movement, any specific historical course for any particular river or
boundary, and no specific identification of any composer with any
later or earlier population. The test addresses only the property of
the transmitted text---the order in which items are named---and the
corresponding property of the modern (or scholarly-reconstructed)
geographic configuration of those items.

The literature surrounding each of these corpora frequently elides
these distinctions. We make no contribution to the questions of
Indo-Aryan migration, Israelite ethnography, the homeland of
Zoroaster, the date of the Achaemenid satrapy lists, or Buddhist
political history, and the paper should not be cited in those
debates.

% =====================================================================
\section{Methods}\label{sec:methods}

\subsection{Geographic identifications}\label{sec:ids}

For each corpus, we adopt the mainstream scholarly identification of
each named entity, taking centroid coordinates from canonical
secondary sources or, for rivers in the Nadīstuti, from
programmatic queries to the OpenStreetMap database. Identifications,
coordinates, confidence levels, and source citations are recorded
per-row in the corpus CSVs deposited with this paper
\citep{repo2026}.

For Ṛgveda~10.75.5 we follow \citet{witzel1995} and
\citet{macdonell-keith1912}, with one philological correction: the
standard Aufrecht text places Asikn\={\i} (Chenab) at position~6 and
Marudv\d{r}dh\={a} at position~7, not the other way around as some
secondary sources have it. We use the standard order.

For the Mahājanapadas we adopt the identifications of
\citet{rhysdavids1903}, \citet{law1973}, and \citet{bronkhorst2007};
the order is the Anguttara Nikāya I.213 sequence
(Kāsi, Kosala, Aṅga, Magadha, Vajji, Malla, Cetiya, Vaṃsa, Kuru,
Pañcāla, Maccha, Sūrasena, Assaka, Avantī, Gandhāra, Kamboja).

For Vendīdād Fargard 1 we adopt \citet{gnoli1980} and
\citet{witzel2000homeland}; eleven of sixteen identifications are
high-confidence, and the five lower-confidence cases
(Airyana~Vaējah, Urvā, Caxra, Varəna, Ra\.{n}gh\={a}) are flagged in the
CSV.

For Genesis~10 we adopt \citet{hess1993} for personal-name analyses
and \citet{westermann1984} and \citet{wenham1987} for the
broader-scholarly identifications. Of seventy-eight named entities
in the canonical Masoretic Text order, fifty are identified with
high or medium confidence; sixteen are either unidentified or
speculative and are dropped from the primary analysis.

\subsection{Test statistics}\label{sec:stats}

\paragraph{Kendall's $\tau$ with permutation null.}
We compute Kendall's $\tau$ \citep{kendall1938}---preferable to
Spearman's $\rho$ for short ordinal sequences, where it carries a
more direct probabilistic reading \citep{newson2002}---between the
sequence-position rank ($1\ldots n$) and the geographic-coordinate
rank of the identified items, on three axes: longitude, latitude,
and the first principal component of the (longitude, latitude)
matrix. A two-sided $p$-value is generated from a 10\,000-iteration
permutation null distribution.

\paragraph{Procrustes superimposition.}
The Kendall-against-longitude and Kendall-against-latitude pair is
not statistically independent when the items lie along a diagonal
sweep, as is common in Eurasian geography. We therefore add a
Procrustes test \citep{procrustes-orig} that superimposes the
sequence-position 1-D embedding onto the 2-D
(longitude, latitude) configuration after optimal rotation,
scaling, and translation; the residual sum of squared distances
$m^2$ is the test statistic, with a permutation null for the
one-sided $p$-value (smaller $m^2$ indicates better alignment).
This addresses the redundant-axis-test concern noted by
\citet{peresnetojackson2001}.

\paragraph{Mantel test.}
Kendall's $\tau$ and Procrustes both assume that the ordering is
\emph{monotonic} with respect to some geographic direction. A list
may, however, be spatially coherent without being monotonic: adjacent
items consistently nearby, but the path through space wandering
rather than directional. The Mantel test \citep{mantel1967} detects
this by correlating the entries of two $n\times n$ pairwise-distance
matrices: $|i-j|$ in the sequence vs.\ the great-circle distance
between items $i$ and $j$. We report the Pearson correlation $r$
between the upper triangles of the two matrices, with a one-sided
permutation null. A small positive $r$ with significant $p$ indicates
that nearby-in-text items tend to be nearby-on-the-ground, even when
the global monotonic ordering is weak.

\paragraph{Constrained permutation null.}
The Ṛgveda corpus is the only one of the five in fixed verse meter.
For it we also compute a permutation null restricted to orderings
that preserve the syllable shape of the verse position by position.
This answers the metrical-determinism objection \citep{arnold1905}:
that in a quantitative-meter text the choice of order may be partly
forced by syllable counts. With shape buckets $\{2:1, 3:3, 4:6\}$,
the meter-admissible set contains
$1!\,\times\,3!\,\times\,6! = 4{,}320$ orderings. The other four
corpora are prose, and the test does not apply.

\paragraph{Bayes factor model comparison.}
Following \citet{kassraftery1995} we compute Bayes factors among
three generative models of the observed ordering:
$M_\text{uniform}$ (sequence drawn uniformly from $n!$
permutations), $M_\text{meter}$ (sequence drawn uniformly from the
meter-admissible set; applicable only to verse corpora), and
$M_\text{geog}(\theta)$ (sequence drawn from a Mallows model
\citep{mallows1957, flignerverducci1986} centered on the
geographically sorted ordering with dispersion $\theta$). Bayes
factors are reported across a grid $\theta \in \{0.5, 1, 1.5, 2, 3\}$.

\paragraph{Reproducibility.}
All analysis code, coordinate CSVs, and the random seed
(\texttt{20260518}) are deposited at \citep{repo2026}.
Dependencies: \texttt{numpy}, \texttt{scipy}, \texttt{pandas},
\texttt{matplotlib}; versions pinned in \texttt{requirements.txt}.

% =====================================================================
\section{Per-corpus results}\label{sec:results}

\subsection{Ṛgveda 10.75.5 (Nadīstuti-sūkta, $n=10$)}\label{sec:rv}

The ten-river enumeration of the Nadīstuti-sūkta yields
$\tau_\text{lon} = -0.911$ with permutation $p < 10^{-4}$
(10\,000 permutations, seed \texttt{20260518}). The Procrustes
superimposition yields $m^2 = 0.179$, $p < 10^{-4}$. The
constrained-permutation null (over the 4\,320 meter-admissible
orderings) finds six orderings yielding $|\tau| \geq 0.911$, for
$p_\text{meter} = 0.0014$: metrical determinism alone cannot
explain the alignment. The Bayes factor against $M_\text{uniform}$
is $\approx 2\times 10^4$ at $\theta = 1.5$; against
$M_\text{meter}$ it is $\approx 24$. The $\tau$ is negative because
longitude falls through the verse, from the eastern-most Ga\.{n}g\={a}
(81.9$^\circ$\,E) to the western-most Su\d{s}om\={a}/Soan
(72.7$^\circ$\,E).

Because RV~10.75.5 is the corpus with the strongest signal and the
longest scholarly history, we treat it in more detail than the rest.
Appendix~\ref{app:rv} gives the ten rivers and their coordinates
(Table~\ref{tab:rivers}), the null distribution
(Figure~\ref{fig:null}), and four sensitivity analyses---alternative
river identifications, paleocourse substitution for the Sarasvatī
and Sutlej, alternative geographic axes, and Gaussian coordinate
perturbation. The result holds in all of them.

As a further check that the result is not an artefact of the
OpenStreetMap coordinate source, we re-derived the centroids of the
nine OSM-traced rivers from MERIT-Hydro, a NASA-SRTM-derived global
hydrography product \citep{yamazaki2019}: for each river we took the
unweighted centroid of all MERIT-Hydro upstream-drainage-area pixels
exceeding a river-network threshold and lying within $0.2^\circ$ of
the OSM polyline. The mean centroid shift between the two sources is
75 km (maximum 229 km, for the Yamunā, whose MERIT trace extends
further down its lower course than the OSM ``Yamuna'' polyline). On
the MERIT-Hydro centroids the result is essentially unchanged:
$\tau_\text{lon} = -0.867$ ($p = 0.0001$), Procrustes $m^2 = 0.227$
($p < 10^{-4}$), Mantel $r = +0.715$ ($p < 10^{-4}$). The geographic
ordering of RV~10.75.5 does not depend on the choice of coordinate
database.

\subsection{Sixteen Mahājanapadas ($n=16$)}\label{sec:mahaj}

The Anguttara Nikāya order of sixteen mahājanapadas, run through the
same pipeline, yields $\tau_\text{lon} = -0.783$, $p < 10^{-4}$, and
Procrustes $m^2 = 0.428$, $p < 10^{-4}$. The Bayes factor against
$M_\text{uniform}$ peaks near $7.7\times 10^4$ at $\theta = 1.5$,
``decisive'' on the Kass-Raftery scale. The canonical order is only
13 inversions, out of a possible 120, away from a perfect
east-to-west sort.

The latitudinal signal is weaker ($\tau_\text{lat} = +0.22$,
$p = 0.27$). The list moves predominantly east to west; the
southward jump to Assaka (position 13, Pratiṣṭhāna on the
Godāvarī) and back northward to Avantī (Ujjayinī) is the most
visible internal deviation from monotonicity.

\subsection{Vendīdād Fargard 1 ($n=16$)}\label{sec:avesta}

The sixteen lands of Fargard 1, in the canonical Old Avestan order
(Airyana Vaējah, Sughdha, Mōuru, B\={a}xd\={\i}, Nisāya, Harōiva,
Vaēkərəta, Urvā, Xnəṇta, Haraxvait\={\i}, Haētumant, Ragh\={a},
Caxra, Varəna, Hapta Həṇdu, Ra\.{n}gh\={a}), yield $\tau_\text{lon} =
-0.117$, $p = 0.56$, and Procrustes $m^2 = 0.975$, $p = 0.67$.
There is a marginal latitudinal signal ($\tau_\text{lat} = -0.38$,
$p = 0.05$), but it does not survive the Procrustes joint test.
Bayes factors for $M_\text{geog}(\theta)$ against
$M_\text{uniform}$ are all $\ll 1$ across the $\theta$ grid: the
Mallows model centered on either-direction longitudinal sort fits
the data \emph{worse} than the uniform prior. The Kendall distance
to the optimal east-to-west sort is 53 of 120, near the random
expectation of 60.

The five lower-confidence identifications (Airyana Vaējah, Urvā,
Caxra, Varəna, Ra\.{n}gh\={a}) collectively cover roughly $35^\circ$ of
longitudinal latitude across published alternative placements, and
we cannot rule out that a better-justified identification of these
five would produce a stronger signal. Within the eleven
high-confidence identifications alone, however, $\tau_\text{lon}$
remains weak ($\tau = -0.21$, $p = 0.35$). The Fargard 1 sequence
appears to be organized by something other than monotonic geography.

\subsection{Genesis 10 (Table of Nations, $n=62$ identified)}\label{sec:gen10}

Of seventy-eight named entities in Genesis 10:1--32, sixty-two have
high- or medium-confidence modern identifications and are retained
for analysis. The flat textual-sequence test yields $\tau_\text{lon} =
+0.174$, $p = 0.05$, and Procrustes $m^2 = 0.91$, $p = 0.005$:
statistically significant on the joint test but with very small
effect size.

But Genesis 10 is not a flat list. It is a genealogy in three
branches---the descendants of Japheth, of Ham, and of Shem---and
within each branch the latitudinal signal is stronger, and points
a different way:

\begin{itemize}
\item Japheth ($n=12$): $\tau_\text{lat} = -0.42$, $p = 0.06$ (weak).
\item Ham ($n=36$): $\tau_\text{lat} = +0.42$, $p < 0.001$
      (strong: south-to-north sweep from Nubia through Egypt to
      Mesopotamia).
\item Shem ($n=14$): $\tau_\text{lat} = -0.56$, $p = 0.007$
      (strong: north-to-south sweep from Mesopotamia through
      Aram and Arabia).
\end{itemize}

The sub-sections cancel at the flat-textual level. Genesis 10
encodes geographic information, but hierarchically rather than as a
single monotonic sweep.

\subsection{Bīsutūn dahyāva ($n=23$)}\label{sec:bisutun}

Column~I of Darius I's Bīsutūn (Behistun) inscription enumerates
twenty-three \textit{dahyāva} ``lands'' subject to Achaemenid
imperial authority. We follow the standard textual order
\citep{kent1953, schmitt2014} (Pārsa, Ūvja, Bābiruš, A\d{t}h\=ur\=a,
Arabāya, Mudrāya, ``those by the sea,'' Sparda, Yauna, Māda,
Armina, Katpatuka, Par\d{t}hava, Zraka, Haraiva, Uvārazmiš, Bāxtriš,
Suguda, Gadāra, Saka, \d{T}hataguš, Harauvatiš, Maka), with nineteen
identifications high-confidence and four (\#5 Arabāya, \#7 Drayahyā,
\#20 Saka, \#21 \d{T}hataguš) marked medium-confidence per
\citet{briant2002}.

The result sits between the Indic corpora and the Avesta:
$\tau_\text{lon} = +0.47$, $p = 0.001$ (moderate west-to-east),
Procrustes $m^2 = 0.51$, $p = 0.0003$, Mantel $r = +0.50$,
$p < 10^{-4}$. The order is not cleanly monotonic. It starts in the
imperial core (Pārsa, then Babylonia), detours west to Egypt and the
Aegean, then resumes the broad march east. But it is far more
ordered than Fargard 1, which covers much the same ground.

Bīsutūn against Vendīdād Fargard 1 is the closest thing to a
controlled comparison in this set of corpora. Both are Iranian, both
of roughly the same date (Bīsutūn \textit{c.}\ 520 BCE; Fargard 1
later in its present form but drawing on older material), and they
name much the same stretch of the Iranian world. Where they part is
genre: Bīsutūn is an administrative declaration of imperial
dominion, Fargard 1 a cosmological account of the lands Ahura~Mazdā
created. Bīsutūn encodes geography; Fargard 1 does not.
Section~\ref{sec:suggests} takes up what to make of that.

% =====================================================================
\begin{figure}[t]
\centering
\includegraphics[width=0.95\linewidth]{fig2_corpus_comparison.pdf}
\caption{Sequence position against longitude, four of the five
corpora shown. Each panel gives the observed scatter, a linear fit
(dashed grey), and the Kendall $\tau$ and Procrustes $m^2$. The two
Indic corpora trace a clean east-to-west line; the Avesta scatters;
Genesis 10 shows only a weak global slope, which §\ref{sec:gen10}
resolves into hierarchical sub-section structure.}
\label{fig:comparison}
\end{figure}

% =====================================================================
\section{Comparative analysis}\label{sec:comparative}

\begin{table}[t]
\caption{Cross-corpus comparison of geographic encoding strength.
$\tau_\text{lon}$: Kendall's $\tau$ vs.\ longitude rank.
$m^2$: Procrustes disparity (smaller = better alignment).
Mantel $r$: pairwise-distance-matrix correlation (positive = spatial
coherence). All $p$-values are two-sided for $\tau$ and one-sided
for Procrustes and Mantel.}
\label{tab:comparison}
\centering\footnotesize
\setlength{\tabcolsep}{4.5pt}
\begin{tabular}{lrrrrrrr}
\toprule
Corpus & $n$ & $\tau_\text{lon}$ & $p_\tau$ & $m^2$ & $p_\text{Proc}$
       & Mantel $r$ & $p_\text{Mantel}$ \\
\midrule
\d{R}gveda 10.75.5          & 10 & $-0.911$ & $<10^{-4}$ & $0.18$ & $<10^{-4}$ & $+0.767$ & $<10^{-4}$ \\
Mah\={a}janapadas (AN I.213)  & 16 & $-0.783$ & $<10^{-4}$ & $0.43$ & $<10^{-4}$ & $+0.617$ & $<10^{-4}$ \\
B\={\i}sut\={u}n dahy\={a}va (DB I.6) & 23 & $+0.472$ & $0.001$    & $0.51$ & $0.0003$   & $+0.505$ & $<10^{-4}$ \\
Vend\={\i}d\={a}d Fargard 1      & 16 & $-0.117$ & $0.56$     & $0.98$ & $0.67$     & $+0.112$ & $0.15$ \\
Genesis 10 (flat)             & 62 & $+0.174$ & $0.05$     & $0.91$ & $0.005$    & $+0.172$ & $0.0004$ \\
\midrule
\multicolumn{8}{l}{\textit{Genesis 10 within-subsection (Kendall $\tau$ vs.\ latitude):}}\\
\quad Japheth (vv. 2--5)      & 12 & $-0.42$  & $0.06$     & ---    & ---       & ---      & ---       \\
\quad Ham (vv. 6--20)         & 36 & $+0.42$  & $<10^{-3}$ & ---    & ---       & ---      & ---       \\
\quad Shem (vv. 21--31)       & 14 & $-0.56$  & $0.007$    & ---    & ---       & ---      & ---       \\
\bottomrule
\end{tabular}
\end{table}

Table~\ref{tab:comparison} collects the cross-corpus numbers. Four
points come out of it.

\paragraph{The Indic lists run in a straight line.}
The Ṛgveda Nadīstuti and the Anguttara Mahājanapadas both sweep
east to west, $|\tau| > 0.78$ on longitude. Out of the 10!
orderings of the RV rivers, and the 16! of the mahājanapadas, the
attested order falls in the top 0.01\,\% closest to a longitudinal
sort.

\paragraph{The Avesta is ordered by something other than geography.}
Fargard 1 covers the same ground as the Indic texts, and the
Iranian plateau besides, but its order is not monotonic on any axis.
That matches \citet{gnoli1980}'s reading that the Fargard 1
sequence follows mythic-cosmographic priority rather than an
itinerary---a qualitative claim to which this is the quantitative
counterpart.

\paragraph{Genesis~10 is ordered hierarchically.}
Read flat, the sequence is barely ordered ($\tau = +0.174$). But the
Mantel test still finds real spatial coherence ($r = +0.172$,
$p = 0.0004$)---neighbouring names are neighbours on the ground even
where the global direction washes out. The three Noachic branches
explain how both can be true at once:

\begin{itemize}
\item Japheth's descendants (positions 1--14) cluster in Anatolia,
      Greece, and the Caucasus. Their internal ordering is weakly
      monotonic against latitude ($\tau_\text{lat} = -0.42$,
      $p = 0.06$).
\item Ham's descendants (15--52) sweep south-to-north from Nubia
      and Egypt through the Levant to Mesopotamia
      ($\tau_\text{lat} = +0.42$, $p < 10^{-3}$).
\item Shem's descendants (53--78) sweep north-to-south from
      Mesopotamia and Aram through Edom to south Arabia
      ($\tau_\text{lat} = -0.56$, $p = 0.007$).
\end{itemize}

The three sub-section signals are individually strong but they
cancel when concatenated, because Ham and Shem run in opposite
latitudinal directions. Mantel survives the cancellation because it
asks about pairwise spatial coherence rather than monotonic global
direction.

This fits the long-standing literary-critical reading
\citep[e.g.,][]{westermann1984} of Genesis~10 as a layered
composite: each Noachic line was a self-contained regional list
before it was hung on the genealogical frame, and the
concatenation kept each list's internal geographic logic while
hiding it at the surface. The quantitative pattern is independent
evidence for that reading---arrived at without any literary-critical
assumption, only by counting.

\paragraph{Length does not explain the pattern.}
The Mahājanapadas ($n=16$) and Fargard 1 ($n=16$) are the same
length but land at opposite ends of the encoding scale. The
Mahājanapadas and the Ṛgveda ($n=10$) are six items apart in length
but score almost the same. Whatever drives the cross-corpus spread
in $\tau$, it is not sample size.

\paragraph{Genre tracks the pattern; language family does not.}
Bīsutūn and Fargard 1 are both Iranian, close in date, and name
overlapping regions, yet one is geographically ordered and the other
is not. The two ordered Indic corpora are a millennium apart and
belong to different genres again---a Vedic hymn and a Buddhist
canonical list---but both clear $|\tau| > 0.78$. Genesis~10's
ordering lives inside its sub-sections, as one would expect of a
genealogy stitched together from older regional lists. Line these
up and the variable that sorts them is genre and compositional
purpose, not the language a text happens to be in or the century it
was fixed. This is a reading the five corpora suggest; with $n=5$ it
cannot be more than that, and §\ref{sec:extensions} sketches the
larger sample that could test it.

\paragraph{Identification uncertainty is bounded.}
Genesis 10's 16 unidentified items and the Avesta's 5
low-confidence items could in principle move the corpus-specific
$\tau$ values. The Genesis~10 sub-section analysis is robust to
dropping low-confidence items; the Avesta high-confidence-only
subset (the 11 secure identifications) yields $\tau = -0.21$,
$p = 0.35$, still null. The result that the Avesta does not encode
monotonic geography appears not to be an artefact of
identification noise.

% =====================================================================
\section{Discussion}\label{sec:discussion}

\subsection{What the pattern suggests}\label{sec:suggests}

The five corpora fall into three kinds of geographic encoding:

\begin{enumerate}
\item \textbf{Strongly monotonic} (RV~10.75.5, the Mahājanapadas).
      The items run along a single east-to-west line, as though the
      composer were reading them off a mental axis. Both corpora are
      Indic; nothing else connects them in date or genre.
\item \textbf{Clustered, with a loose global direction}
      (Bīsutūn dahyāva). Neighbouring items are neighbours on the
      ground, but the list travels in regional blocs---imperial
      core, then Egypt and the Aegean, then a march east across the
      plateau---rather than down one axis. That is what listing
      satrapies in administrative-priority order would look like.
\item \textbf{Hierarchical} (Genesis~10). Each Noachic branch is
      latitudinally ordered on its own; the three branches point
      different ways and cancel once concatenated. A genealogy
      assembled from older regional lists would behave exactly so.
\end{enumerate}

Fargard 1 fits none of the three. It returns a null on every test.
It is also the one corpus here that is plainly cosmological---Ahura
Mazdā's creation of the lands, set line by line against Aŋra
Mainyu's counter-creation of evils---and a list built around
creational priority has no reason to be geographic. Set beside
Bīsutūn, which names much the same world in the same century and
\emph{is} ordered, Fargard 1 forms the cleanest natural experiment
the five corpora contain: hold language family and date fixed, vary
genre, and the geographic signal appears or disappears with it.

Why genre should do this, we cannot say from five texts. Three
factors are worth pulling apart in a larger sample: genre proper
(an invoked hymn, a canonical political list, an imperial
declaration, a cosmology, a genealogy); date, and the depth of oral
transmission behind the written text; and how linear the composer's
own mental map was likely to be---someone whose world runs along a
single dominant axis may simply default to it.

\subsection{What it does not imply}\label{sec:nogo}

A result here says nothing about whether any tradition had special
geographic-cognitive access, nothing about the identity, ethnicity,
or migration history of any composer, and nothing for or against any
particular philological identification beyond what published
scholarship already supports. The tests bear only on whether the
listed order correlates with the modern coordinates of the
identified items.

\subsection{Limitations}\label{sec:limits}

The most serious limitation is the small number of corpora. Five
texts differing at once in length, identification quality, genre,
and date leave those factors hopelessly confounded; the genre
reading of §\ref{sec:suggests} is a hypothesis the present data
suggest, not one they establish. Isolating any single factor would
need a sample several times larger.

Two measurement choices are also worth flagging. A centroid is a
coarse summary of a long linear feature, whether a river or a
province; for the Mahājanapadas and the broader Genesis~10 regions,
choosing the capital city instead of the administrative-area mean
shifts individual coordinates by tens of kilometres. We checked that
the cross-corpus comparison survives that choice, but it is a
commitment, not a fact of the data. And the Avesta and Genesis~10
corpora carry a non-trivial share of lower-confidence
identifications; dropping them leaves the qualitative conclusions in
place, but a full robustness study of the identifications is a
separate piece of work.

Finally, the Mallows model behind the Bayes factor is one defensible
choice among several. A Cayley- or Hamming-distance Mallows, or a
non-Mallows prior, could move the numbers. We use the
Kendall-distance form because it is the one that matches Kendall's
$\tau$ as the test statistic.

\subsection{Extensions}\label{sec:extensions}

The same method runs on any fixed-order enumeration with
identifiable items. Obvious candidates for a larger study are the
Mahābhārata Bhīṣmaparvan river list, Strabo's enumeration of Asian
peoples in books 11--15, the directional mountain sequences of the
Shan Hai Jing, and the legendary regions of Avesta Yašt 19. A
comparative set of twenty or thirty such lists would begin to let
the genre, length, and date factors be separated. Two refinements of
the method itself are also worth pursuing: distance-matrix
correlations beyond the single Mantel statistic used here, and
constrained nulls whose prior is built from attested intertextual
list-traditions rather than from meter alone.

% =====================================================================
\section{Conclusion}\label{sec:conclusion}

Geographic ordering is not a universal property of premodern lists.
Of the five enumerations tested, the Ṛgveda Nadīstuti and the
Anguttara Mahājanapadas trace a strong monotonic east-to-west line
($|\tau| > 0.78$, $p < 10^{-4}$); the Bīsutūn dahyāva is ordered
more loosely, by satrapal bloc; Genesis 10 is ordered within its
sub-sections but not across them; and Vendīdād Fargard 1 shows no
geographic order at all. The cleanest signal in the data is that
genre tracks these outcomes better than language family does. We
claim nothing past the comparisons themselves, and the method
carries over to any other fixed-order enumeration.

% =====================================================================
\section*{Acknowledgements}
This analysis depended on the open data of the OpenStreetMap
project (Overpass API), the publicly-available scholarly
identifications in the works cited, and the canonical electronic
texts maintained by GRETIL, TITUS, Sefaria, and the Chinese Text
Project. No funding sources are reported.

\section*{Data and code availability}
All data, code, cached Overpass responses, and the random seed
(\texttt{20260518}) used for the permutation tests are available at
\citep{repo2026}. Each corpus's coordinate CSV, including
identification provenance and confidence levels, is provided in
\texttt{corpora/}. Reproduction instructions in
\texttt{REPRODUCING.md}.

% =====================================================================
\appendix
\section{The R\d{g}veda 10.75.5 corpus in detail}\label{app:rv}

RV~10.75.5 is the corpus with the strongest geographic signal and
the longest scholarly history---it has been read as evidence for
Vedic-period geography since \citet{roth1846} and \citet{zimmer1879}.
Section~\ref{sec:rv} treats it more briefly than that history
warrants. This appendix gives the underlying data and the
sensitivity analyses behind the one-line claim there that ``the
result holds.''

\paragraph{The verse.}
RV~10.75.5 names its ten rivers in a single \d{r}c, here in the
Aufrecht text \citep{aufrecht1877}, confirmed against
\citet{vannooten-holland1994}:

\begin{quote}
\textit{ima\d{m} me ga\.{n}ge yamune sarasvati\\
\'sutudri stoma\d{m} sacat\={a} paru\d{s}\d{n}y\={a}\\
asikny\={a} marudv\d{r}dhe vitastay\={a}\\
\={a}rj\={\i}k\={\i}ye \'s\d{r}\d{n}uhy\={a} su\d{s}omay\={a}}
\hspace{1.5em}(RV 10.75.5)
\end{quote}

One transcription point is worth flagging, because some secondary
sources have it wrong: the standard text places Asikn\={\i} (Chenab)
at position~6 and Marudv\d{r}dh\={a} at position~7, directly after
Paru\d{s}\d{n}\={\i} (Ravi)---not the reverse. The analysis uses the
standard order.

\paragraph{Identifications and coordinates.}
Table~\ref{tab:rivers} lists the ten rivers, their mainstream
identifications \citep{witzel1995, macdonell-keith1912}, and the
centroid coordinates used. The two residual disagreements in the
literature---the referent of \={A}rj\={\i}k\={\i}y\={a} and the
historical course of the Sarasvat\={\i}---are well characterised
\citep{witzel1999, brereton1981} and are carried through the
sensitivity analysis. Nine centroids were extracted from
OpenStreetMap; Marudv\d{r}dh\={a}, for which no OSM-tagged river of
that name exists, is hand-curated to the Chenab--Jhelum doab and
exercised in the sensitivity analysis below.

\begin{table}[t]
\caption{The ten rivers of RV~10.75.5 in verse order, with
mainstream identifications and centroid coordinates.
``$n_\text{nodes}$'' is the number of OSM geometry nodes behind the
centroid; Marudv\d{r}dh\={a} is hand-curated.}
\label{tab:rivers}
\centering\small
\begin{tabular}{rllrrr}
\toprule
Pos. & Vedic name & Identified river & Lon (\textdegree\,E) &
Lat (\textdegree\,N) & $n_\text{nodes}$ \\
\midrule
1  & Ga\.{n}g\={a}      & Ganges               & 81.88 & 26.87 & 15\,159 \\
2  & Yamun\={a}         & Yamuna               & 78.09 & 29.10 & 19\,337 \\
3  & Sarasvat\={\i}     & Ghaggar--Hakra       & 75.76 & 29.87 & 12\,673 \\
4  & \'Sutudr\={\i}     & Sutlej               & 77.09 & 31.30 &  4\,047 \\
5  & Paru\d{s}\d{n}\={\i} & Ravi               & 74.22 & 31.48 & 12\,047 \\
6  & Asikn\={\i}        & Chenab               & 73.68 & 31.97 & 12\,824 \\
7  & Marudv\d{r}dh\={a} & Chenab--Jhelum doab (uncertain) & 73.10 & 32.20 & --- \\
8  & Vitast\={a}        & Jhelum               & 73.51 & 33.27 & 13\,624 \\
9  & \={A}rj\={\i}k\={\i}y\={a} & Haro          & 72.84 & 33.86 &  2\,018 \\
10 & Su\d{s}om\={a}     & Soan (Sohan/Sawan)   & 72.75 & 33.37 &  3\,190 \\
\bottomrule
\end{tabular}
\end{table}

\paragraph{Primary test.}
Kendall's $\tau$ between verse order and longitude is
$\tau = -0.911$. Across 10\,000 permutations no permuted statistic
reached $|\tau|$, so $p < 10^{-4}$. The null distribution
(Figure~\ref{fig:null}) is near-symmetric about zero, mean
$-0.001$, standard deviation $0.251$, central 95\,\% interval
$[-0.511, +0.467]$.

\begin{figure}[t]
\centering
\includegraphics[width=0.74\linewidth]{fig1_null_distribution.pdf}
\caption{Null distribution of Kendall's $\tau$ between verse-order
rank and longitude, over 10\,000 random permutations of the
verse-order labels. The observed $\tau = -0.911$ falls far outside
the central 95\,\% interval of the null.}
\label{fig:null}
\end{figure}

\paragraph{Sensitivity analyses.}
Four checks, summarised in Table~\ref{tab:sensitivity}.
\textit{Geographic axis:} $\tau$ against latitude is $+0.956$
($p<10^{-4}$) and against the first principal component $+0.911$
($p=0.0002$)---the latitude figure is not independent evidence,
since the rivers lie on a NW--SE diagonal, but it rules out an
axis-choice artefact. \textit{Paleocourse:} substituting published
paleocourse centroids---the Ghaggar--Hakra reconstruction of
\citet{giosan2012} and \citet{clift2012} for the Sarasvat\={\i},
whose precise course is still under investigation
\citep{singh2017}, and the older pre-capture course of
\citet{yashpal1980} for the \'Sutudr\={\i}---gives
$\tau = -0.778$ ($p = 0.001$): attenuated, still strong.
\textit{Alternative identifications:} across six variant
identifications of the disputed \={A}rj\={\i}k\={\i}y\={a} and
Marudv\d{r}dh\={a}, $\tau$ stays between $-0.689$ and $-0.911$, all
$p \leq 0.005$; dropping the contested Marudv\d{r}dh\={a} entirely
($n=9$) gives $\tau = -0.944$. \textit{Coordinate perturbation:}
Gaussian noise at $\sigma = 0.5^\circ$, 1\,000 iterations, yields a
95\,\% interval for $\tau$ of $[-0.956, -0.644]$, median $-0.822$.

\begin{table}[t]
\caption{RV~10.75.5 sensitivity-analysis summary. ``$p$ two-sided''
is the share of permuted $|\tau|$ values $\geq$ the observed, from
10\,000 permutations.}
\label{tab:sensitivity}
\centering\footnotesize
\setlength{\tabcolsep}{5pt}
\resizebox{\textwidth}{!}{\begin{tabular}{lrrl}
\toprule
Analysis & $\tau$ & $p$ (two-sided) & Note \\
\midrule
Primary: longitude & $-0.911$ & $2.98\times10^{-5}$ & modern centroids \\
Axis: latitude & $+0.956$ & $5.51\times10^{-6}$ & rivers also trend N-S \\
Axis: PC1 & $+0.911$ & $2.98\times10^{-5}$ & principal component of $(\mathrm{lon},\mathrm{lat})$ \\
Paleocourse substitution & $-0.778$ & $9.46\times10^{-4}$ & paleo-Sarasvat\={\i}, paleo-Sutlej \\
ID: \={A}rj.=Haro; Mar.=Chenab–Jhelum doab & $-0.911$ & $2.98\times10^{-5}$ &  \\
ID: \={A}rj.=Haro; Mar.=eastern alternative & $-0.911$ & $2.98\times10^{-5}$ &  \\
ID: \={A}rj.=Soan; Mar.=Chenab–Jhelum doab & $-0.689$ & $0.0047$ &  \\
ID: \={A}rj.=Soan; Mar.=eastern alternative & $-0.689$ & $0.0047$ &  \\
ID: \={A}rj.=upper Indus tributary; Mar.=Chenab–Jhelum doab & $-0.867$ & $1.15\times10^{-4}$ &  \\
ID: \={A}rj.=upper Indus tributary; Mar.=eastern alternative & $-0.867$ & $1.15\times10^{-4}$ &  \\
ID: \={A}rj.=Haro; Mar.=DROPPED & $-0.944$ & $4.96\times10^{-5}$ &  \\
Coord. perturb. ($\sigma=0.5^\circ$, n=1000) & $-0.822$ & --- & 95\% CI [-0.956, -0.644] \\
\bottomrule
\end{tabular}
}
\end{table}

Every variant leaves the result in the same place: the order of
rivers in RV~10.75.5 is strongly, and not accidentally, aligned
with longitude.

\bibliographystyle{plainnat}
\bibliography{references}

\end{document}
